\chapter*{Abstrakt}
\addcontentsline{toc}{chapter}{Abstrakt}

Tato diplomová práce se zabývá marketingovým výzkumem zaměřeným na posouzení potenciálu vstupu nové značky mraženého ovoce v~čokoládě na český trh. Práce vychází z~identifikace mezery na trhu, kde v~současnosti dominuje zahraniční značka Franui bez přímé lokální konkurence. Cílem výzkumu je analyzovat spotřebitelské preference, cenovou citlivost a~vnímání kategorie mraženého ovoce v~čokoládě českými spotřebiteli, a~tím poskytnout podklady pro rozhodnutí o~uvedení nové značky Berrie na trh.

Výzkum využívá smíšenou metodologii. V~kvalitativní fázi jsou realizovány dvě skupinové diskuse (focus groups) zaměřené na generaci~Z a~matky s~dětmi, doplněné o~senzorickou analýzu formou Central Location Testu. Kvantitativní fáze zahrnuje online dotazníkové šetření (\gls{cawi}) na vzorku 400 respondentů s~využitím metody Price Sensitivity Meter pro stanovení optimální cenové hladiny.

Výsledky práce poskytují komplexní pohled na potenciál nového produktu a~formulují doporučení pro marketingovou strategii značky Berrie při vstupu na český trh.

\vspace{5mm}
\noindent\textbf{Klíčová slova:} marketingový výzkum, mražené ovoce v~čokoládě, vstup na trh, spotřebitelské preference, focus groups, senzorická analýza, cenová citlivost, Franui, Česká republika

% ============================================================================
\newpage
\chapter*{Abstract}
\addcontentsline{toc}{chapter}{Abstract}

This diploma thesis presents a~marketing research study aimed at assessing the potential for introducing a~new brand of chocolate-covered frozen fruit to the Czech market. The research stems from the identification of a~market gap, where the foreign brand Franui currently dominates without a~direct local competitor. The objective is to analyse consumer preferences, price sensitivity, and the perception of the chocolate-covered frozen fruit category among Czech consumers, thereby providing a~foundation for the market entry decision of the new brand Berrie.

The research employs a~mixed-method approach. The qualitative phase consists of two focus group discussions targeting Generation~Z consumers and mothers with children, complemented by a~sensory evaluation conducted as a~Central Location Test. The quantitative phase comprises an online survey (\gls{cawi}) with a~sample of 400 respondents, utilising the Price Sensitivity Meter method to determine the optimal price point.

The findings offer a~comprehensive assessment of the new product's potential and formulate recommendations for the marketing strategy of the Berrie brand upon entering the Czech market.

\vspace{5mm}
\noindent\textbf{Keywords:} marketing research, chocolate-covered frozen fruit, market entry, consumer preferences, focus groups, sensory analysis, price sensitivity, Franui, Czech Republic

\newpage
